\section{مجموعه‌ای که پروژه‌به‌پروژه رشد می‌کند}
این مخزن را برای استفادهٔ دوباره از درایورهای سیستم‌های نهفته و آموزش روش ساخت آن‌ها نگه می‌دارم. هر قطعه یک پروژهٔ کامل دارد که می‌توان آن را به پوشه‌ای دیگر کپی کرد. خود کتابخانهٔ قابل‌حمل نیز به‌تنهایی قابل انتقال است؛ برای اتصال آن به میکروکنترلر دیگر یک آداپتر سخت‌افزاری کوچک نوشته می‌شود.

گزارش، مستند مشترک مجموعه است و وابستگی مشترک ساخت firmware نیست. یک نسخهٔ انگلیسی و یک نسخهٔ فارسی داریم. هر درایور جدید یک فصل دارد که مدار، کد، روش تست و وضعیت اعتبارسنجی همان پروژه را توضیح می‌دهد. README هر پروژه باید برای استفادهٔ مستقل و بدون این گزارش کافی باشد.

\subsection{فهرست فعلی}
\begin{tabular}{p{27mm}p{57mm}p{49mm}}
\toprule
پروژه & کاربرد & شواهد فعلی\\
\midrule
\lr{DS18B20} & خواندن دمای یک سنسور روی
\lr{STM32F107VCT6 / PB10}
& ساخت و تست نرم‌افزاری؛ تست برد باقی است\\
قالب پایه & شروع مستقل برای پروژهٔ جدید & وضعیت محلی قالب بررسی شود\\
\bottomrule
\end{tabular}

\subsection{جای هر بخش در پروژه}
\begin{LTR}
\begin{lstlisting}[language={},numbers=none]
ds18b20/
  DS18B20.ioc          # hardware configuration
  Core/               # CubeMX-generated initialization
  Drivers/            # local ST HAL and CMSIS
  Application/        # board adapter and hardware demo
  Libraries/OneWire/  # portable wire protocol
  Libraries/DS18B20/  # portable temperature operation
  MDK-ARM/            # Keil project and startup
  platformio.ini      # GCC build and probe selection
  tests/              # simulated portable C tests
  tools/              # source synchronization/checks
  README.md           # standalone usage instructions
  PROJECT_STATUS.md   # evidence and unfinished work
docs/
  english/            # one complete report edition
  persian/            # matching report edition
  tools/              # report build and PDF verification
\end{lstlisting}
\end{LTR}
بخش
\lr{Core}
میکروکنترلر را راه‌اندازی می‌کند. بخش
\lr{Application}
رجیستر GPIO و کلاک را می‌شناسد. بخش
\lr{Libraries}
کد C قابل استفادهٔ دوباره برای صحبت با سنسور را نگه می‌دارد. با این جداسازی، کد دما هم در شبیه‌سازی و هم روی برد دیگر قابل استفاده است.

\subsection{روش مطالعه و بازتولید}
برای اولین نتیجه، مرحله‌های ۱ تا ۴ فصل سنسور را به ترتیب انجام دهید. مرحلهٔ ۵ دو تابع لازم برای پروژهٔ دیگر را آموزش می‌دهد و مراحل ۶ و ۷ نمونه و تغییر پایه را توضیح می‌دهند. مراحل ۸ تا ۱۰ پس از نتیجۀ عملی به پیاده‌سازی می‌پردازند؛ می‌توانید بعداً به آن‌ها برگردید. ساخت موفق، تست شبیه‌سازی و اندازه‌گیری واقعی سه نتیجهٔ جدا هستند. مرحلهٔ نهایی، ثبت تست پذیرش برد است.

PDF دیتاشیت‌ها فقط مرجع محلی‌اند و وارد گیت نمی‌شوند. README و منابع گزارش، لینک رسمی و نسخهٔ منبع را حفظ می‌کنند. PDF نهایی گزارش و سورس قابل ویرایش در گیت نگهداری می‌شوند؛ فونت دانلودشده، تصویر بررسی و خروجی ساخت وارد گیت نمی‌شوند.
